\section{Discussion and Limitations} \label{sec:discussion-and-limitations}
  % Page Policy
  In this work, an issue was observed when using the Knock-Knock tool~\cite{toolKnockKnock} to 
  reverse engineer the \gls{dram} address mapping. Specifically, the tool reported the 
  failure condition “\textit{Valley between peaks not pronounced enough}”. As a 
  consequence, no part of the mapping could be retrieved from the measurements.
  \\
  A closely related issue has been reported in prior work. The authors of~\cite{paperKnockKnock} 
  encountered the same problem on \gls{lpddr}-based devices, including the Raspberry 
  Pi 4 with 4 GB of RAM and the Nvidia Jetson Nano. They attributed this 
  behavior to the presence of a closed page policy, as their latency measurements 
  showed only a single distribution rather than clearly separable peaks. Similar 
  behavior was observed on additional platforms, including the Raspberry Pi 4 
  with both 4 GB and 8 GB of RAM, as well as the newer Raspberry Pi 5 with 4~GB, 
  8~GB, and 16~GB of RAM, when analyzed using the Knock-Knock tool~\cite{toolKnockKnock}.
  \\
  Under a closed page policy, the memory controller immediately closes the 
  currently active row after each memory access~\cite{thesisPagePolicies, paperKnockKnock}. In contrast, an open 
  page policy, which is the default in most systems, keeps a \gls{dram} row active 
  in the row buffer until a subsequent access targets a different row, at which 
  point the currently open row must be closed~\cite{paperAnotherFlip, thesisPagePolicies}.
  \\
  Based on the measurements conducted in this work, the most likely 
  explanation for the observed behavior is not the presence of a closed page 
  policy, but rather an unclear limitation of the Knock-Knock methodology on 
  the evaluated platforms. Although the tool fails with the error message, the measurements 
  demonstrate that row hits and row conflicts can, in principle, be 
  distinguished, indicating that the row buffer is not always closed. This 
  apparent contradiction is further supported by other studies~\cite{paperDRAMMaUT, paperMemoryAware}, which 
  evaluated the same device, namely the Raspberry Pi 4. Consequently, this 
  limitation does not arise from the \gls{dram} configuration itself, but from the 
  constraints of the Knock-Knock methodology. While the underlying memory behavior 
  allows distinguishing row hits from row conflicts, the tool fails to recover the 
  mapping due to its thresholding and peak-detection assumptions.

  % Refresh Control
  After successfully porting TRRespass to ARM, extensive testing was performed 
  over several hours using different configurations and multiple runs in an attempt 
  to bypass \gls{trr} mechanisms and induce bitflips. Despite these efforts, no effective 
  hammering pattern emerged when using the fuzzer.
  \\
  Similar limitations have been reported in~\cite{paperBlacksmith} on some of their tested Micron
  LPDDR4x devices as they were unable to trigger bitflips. Regardless, the official 
  publication does not specify the exact models or their architectures. Especially 
  newer modules manufactured in 2020 were affacted due to improved mitigation 
  mechanisms such as strengthened \gls{trr} implementations. As a result, its effectiveness 
  on \gls{lpddr} devices was shown to be platform-dependent and not guaranteed across all 
  chips.
  \\
  Blacksmith also did not work on these devices out of the box without reverse 
  engineering the \gls{trr} behavior. Its default configurations assumed \gls{trr} distances, 
  hammering amplitudes, and refresh alignments that did not correspond to the actual 
  \gls{trr} behavior of certain LPDDR4x devices. The failure on some LPDDR4x modules was 
  attributed to several factors. First, \gls{trr} distances were larger and more irregular 
  than assumed by the default parameter ranges. Second, effective hammering patterns 
  required longer consecutive hammering phases and precise alignment with specific 
  refresh commands. Third, correct refresh alignment was critical but not accounted 
  for in the initial implementation~\cite{toolBlacksmith}. Finally, the required reverse engineering 
  depended on control over refresh commands, which is generally unavailable outside 
  controlled test platforms.
  \\
  These findings directly relate to the limitations of the present work. In the 
  evaluation, the authors had full control over the refresh commands on 
  their test devices, which enabled them to reverse engineer the \gls{trr} behavior. In 
  contrast, the experimental setup used in this work does not provide access to 
  refresh control, thereby preventing a comparable analysis of \gls{trr} mechanisms. 
  Consequently, the lack of refresh awareness and the restricted fuzzing capabilities 
  limit the reliable identification of effective attack patterns.
  \\
  To make Blacksmith-like approaches work on similar \gls{lpddr} devices without refresh 
  control, alternative strategies are needed. These include probabilistic or exhaustive 
  fuzzing with extended parameter ranges to compensate for unknown \gls{trr} distances and 
  refresh behavior, as well as accepting longer runtimes and repeated attempts, since 
  effective patterns require alignment with refresh intervals that cannot be directly 
  controlled. However, without explicit control over refresh commands, the probability 
  of accidentally achieving the correct refresh alignment remains small, particularly
  given larger \gls{trr} distances and irregular refresh behavior. Achieving such an alignment 
  would require very long runtimes and many repeated attempts, rendering the approach 
  impractical in most realistic settings.
